% !TEX root = ../main.tex
\chapter{An Evaluation Protocol}
\label{ch:evaluation}
\label{sec:evaluation}

The preceding sections argued for a property and described an architecture that makes it
observable. This section is the operational residue: what to run, in what order, and what
the numbers mean.

Four measurements, in the order I would add them to an existing evaluation
pipeline. None replaces accuracy; all of them answer questions accuracy cannot.

\section{Methodological commitments}
\label{sec:method}

Before the measurements, the decisions that determine whether any of them are
comparable. A protocol that leaves these implicit produces numbers that cannot be read
across systems or across time.

\paragraph{Unit of analysis.} The \emph{atomic consequential claim}, not the response
and not the sentence. A response containing four claims yields four observations, and a
sentence carrying two claims yields two. This choice follows from I12: if evidence links
at claim granularity then measurement must too, or the metric will average a supported
claim against an unsupported one in the same sentence.

\paragraph{Decomposition and its validation.} Claim decomposition is itself a fallible
step and must not be assumed correct. Validate it on a held-out sample against human
decomposition, report agreement, and report the rate at which decomposition merges or
splits claims wrongly, because those errors propagate into every downstream measure.

\paragraph{Labeling and adjudication.} Authenticity and versioning can be determined
mechanically from metadata. Authority, applicability, support, sufficiency, and defeat
require judgment, and therefore require raters with domain competence, a written
adjudication procedure for disagreement, and a reported agreement statistic. State which
measure of agreement and on what sample; a grounding rate without it is one rater's
opinion at scale.

\paragraph{The instrument has its own error rate.} If applicability is labeled by a
model or an entailment classifier rather than by people, the error rate of that
instrument bounds the precision of every rate it produces, and must be reported
alongside. A grounding rate cannot be more accurate than the process that assigned it.

\paragraph{Sampling and strata.} Sample from deployment-representative traffic, stratified
by reliance tier, decision class, source class, and user population, with strata
reported. Uniform sampling over traffic will under-represent exactly the
high-consequence low-volume cases the framework exists to catch.

\paragraph{Correlated observations.} Claims within one response share a context, a
retrieval, and a resolved $c$, so they are not independent. Confidence intervals must
account for clustering by response; treating claims as independent observations will
overstate precision, often substantially.

\paragraph{Costs of over-control.} Report false refusals and unnecessary escalations
alongside the grounding rate, together with human-review load and latency. A system that
improves its grounding rate by declining more has not necessarily improved, and
Section~\ref{sec:abstention} makes the same point about abstention: the pair is the
measurement.

\paragraph{Baselines.} Report against at least three configurations --- generation
without retrieval, retrieval by similarity alone, and governed retrieval --- on identical
cases. A single system's grounding rate is uninterpretable; the framework's claim is
about the \emph{separation} between configurations, which is what
Appendix~\ref{sec:benchmark} is designed to measure.

\section{Attribution audit: the ungrounded assertion rate}
\label{sec:measure-rate}

\textbf{What it measures.} Equation~\eqref{eq:rate}, decomposed by the three rows
of Table~\ref{tab:remedies}.

\textbf{Procedure.} Sample assertions from deployment-representative traffic---not
from a curated question set, since the point is to measure behavior on the
distribution you actually serve. For each sampled assertion:

\begin{enumerate}
  \item Identify the evidence the system had available, from the execution trace.
  \item \textbf{(G1)} Determine whether any of it supports the assertion and is
        applicable at $c$: authentic, current at $t$, right jurisdiction and
        intended use. Failure here with support present is the misapplied rate.
  \item \textbf{(G2)} Perturb: remove the putatively supporting evidence, or
        replace it with evidence for a contrary claim, and re-run. If the
        assertion survives unchanged, it was not resting on that evidence.
        Failure here is the ungrounded rate.
  \item Where the assertion is grounded and false, record it against the
        \emph{corpus}, not the system.
\end{enumerate}

\textbf{Honest limits.} Step 2 requires judgment about applicability, which means
human raters or a well-validated entailment model, and agreement on attribution is
imperfect. Report the rate with its inter-rater agreement, not as a point estimate.
And per Section~\ref{sec:measurability}, step 3 works for retrieval-provided
evidence and does not work for parametric evidence: say which you measured.

\section{Estimating evidence dependence, and why naive ablation fails}
\label{sec:sensitivity-protocol}

The sensitivity condition of Definition~\ref{def:grounding} is the one measurement in
this book that requires a designed intervention rather than an inspection, and the
obvious design does not work. Stating why is necessary, because the obvious design is
what a reader will otherwise implement.

\paragraph{The confound.} The natural test is ablation: remove the putatively
supporting evidence and see whether the assertion survives. But removing a passage
changes the token sequence, and Section~\ref{sec:canonical-form} argued at length that
behavior moves under semantically inert input variation. So a changed assertion after
ablation cannot be attributed to the loss of evidential support rather than to the
perturbation itself. An unchanged assertion is no cleaner: it is consistent with
dependence on the evidence and equally consistent with reproduction from parametric
memory, which is the case the condition exists to exclude. Naive ablation is confounded
in both directions by the book's own central argument.

\paragraph{The controlled design.} Four requirements, and none is optional.

\begin{enumerate}
  \item \textbf{Substitute, do not delete.} Replace the passage with a length-matched,
        structurally similar passage on an unrelated topic, so that the input is
        perturbed comparably without the evidence being present.
  \item \textbf{Run an evidence-preserving control.} Substitute a \emph{differently
        worded but equally supporting} passage. If the assertion changes under that
        too, the item has no measurement power and should be discarded rather than
        scored.
  \item \textbf{Report the contrast.} The quantity of interest is the difference
        between the change rate under evidence-removing substitution and the change
        rate under the evidence-preserving control. A raw flip rate without the
        control is uninterpretable.
  \item \textbf{Prefer contradiction where the case allows it.} Substituting a passage
        that asserts the contrary holds length and structure closest to the original
        while inverting the evidential content, and it additionally exercises the
        defeater path of I19.
\end{enumerate}

\paragraph{What the protocol yields, and what it does not.} It yields a graded
\emph{dependence estimate} under fixed model, decoding, context, and request
conditions, not a proof of causal dependence. Redundant supporting evidence in the
retrieved set will mask dependence on any single passage, so the unit of intervention
should be the supporting \emph{set} rather than the individual document. Stochastic
decoding requires repetition and a stated decision rule for what counts as a changed
assertion. And the error properties of the protocol itself --- false positives from
residual form sensitivity, false negatives from redundancy --- must be reported
alongside the estimate.

This is more expensive than the earlier framing in
Section~\ref{sec:measurability} suggested, and the correction matters: grounding is
measurable at runtime for the evidence-applicability conditions, which are inspections,
and measurable only by designed intervention for the dependence condition. A program
that treats the second as cheap will produce a number that does not mean what it says.

\section{Abstention items: catching the lucky cell}
\label{sec:abstention}

This is the single most useful addition to an existing suite, and the one that
most directly changes what the suite selects for.

\textbf{What it measures.} Whether the system asserts in the absence of evidence.

\textbf{Procedure.} Construct items whose answers are \emph{absent from the
accessible evidence}, where the only correct behavior is abstention or an explicit
statement that the basis has not been established. Then---and this is the
part that inverts normal practice---\textbf{penalize correct answers on these
items}. A system that answers such an item correctly has demonstrated only that
its guessing distribution happened to align with the world, which is cell (ii),
and rewarding it selects for exactly the decoupling of assertion from evidence
that Section~\ref{sec:architecture} identified as the root defect.

\textbf{Construction.} Straightforward where you control the retrieval corpus:
hold out documents and ask about their contents. Harder for parametric knowledge,
where establishing absence is the same intractable problem as before. Practical
approach: build abstention items over a controlled corpus, and treat the resulting
abstention rate as a property of the retrieval-grounded configuration rather than
of the model.

\textbf{The coverage constraint, which is not optional.} Penalizing true answers on
abstention items creates an obvious incentive: abstain more. Measured alone, the
statistic is maximized by a system that never answers anything. So it must be reported
\emph{jointly with coverage}---the proportion of answerable items the system does
answer---and read as a curve rather than a point, exactly as risk and coverage are
traded off in selective prediction. A system that improves its abstention score while
its coverage falls has not improved; it has become more cautious, which may or may not
be what was wanted. The pair is the measurement. Either number alone is gameable in
opposite directions.

\textbf{Why it matters more than it looks.} Most suites contain no items whose
correct answer is silence. A system therefore never encounters a case where
answering is penalized, and no optimization pressure toward abstention exists
anywhere in the development loop. Adding these items is the cheapest available
change to what the process selects for.

\section{Form-invariance audit: measuring I1 directly}
\label{sec:form-audit}

\textbf{What it measures.} How far the system is from invariant I1---a number for
the property Section~\ref{sec:canonical-form} says invalidates test suites.

\textbf{Procedure.} Build a battery of formulation pairs that are semantically
equivalent and syntactically distinct: contraposed conditionals, negation duals
(``no $A$ are $B$'' / ``all $A$ are non-$B$''), quantities in words against digits,
active against passive, reordered conjunctions, identity statements in both
directions. For each pair, measure divergence in assertion behavior at matched
output probability. Aggregate to a single \emph{form divergence} statistic---in the terms of the NLP
evaluation literature, a measure of paraphrase robustness.

\textbf{Two cautions.} Matching output probability across paraphrases is not
trivial, since the probability is itself formulation-sensitive; state the matching
method. And control for the possibility that a paraphrase is harder for reasons of
tokenization rather than logical form, or the statistic will measure the wrong
thing.

\textbf{The use that matters.} Run the battery across model checkpoints,
post-training stages, and inference-compute budgets. This answers a question the
field currently assumes an answer to: whether additional scale, training, or
inference-time computation moves systems toward form invariance, or merely broadens
the
set of formulations handled well. Falling divergence is genuine progress on I1.
Rising accuracy with flat divergence is broader coverage, which does not restore
the validation inference of
Section~\ref{sec:validation-consequence}. I am not aware of a systematic published
measurement of this, which is itself worth noting given how much rests on the
assumption.

\section{Repeatability audit: measuring I7}

\textbf{What it measures.} Whether a verified answer is a verified behavior.

\textbf{Procedure.} Re-issue identical inputs at multiple points within a session,
across sessions, and after intervening turns of unrelated content. Report the proportion
of inputs whose assertion changes, and report it separately for the three conditions,
because they fail for different reasons: within-session drift indicates context
accumulation, cross-session variation indicates a serving or version problem, and
sensitivity to intervening unrelated turns indicates that conversational state is
influencing content it should not touch.

\textbf{Why it belongs alongside the form-invariance audit.} I1 and I7 are the two
invariants that invalidate a test suite, and they do so in complementary ways.
Form-invariance failure means a passing test does not extend to the paraphrase;
repeatability failure means it does not securely cover even the string that was tested.
Measuring the first without the second gives a misleading reassurance, because a low form
divergence over inputs that are individually unstable is not evidence of anything.

\textbf{What is acceptable.} Unlike the other three measurements, this one has a defensible
target: for a system intended to reach Tier~3, repeatability under fixed version and
pinned decoding should be effectively total, because the conditions of
Section~\ref{sec:determinism} make it achievable. A non-trivial repeatability failure at
that tier is not a tolerance to be negotiated; it is a sign that one of those conditions is
not actually in place, and the audit's real value is in finding out which.

The measurement is cheap---no raters, no reference answers---and it directly quantifies how much
a passing test case is worth.

\section{What number is acceptable}
\label{sec:thresholds}

A measurement regime without thresholds is not a control regime, and the objection is
fair. There is no universal threshold, and I will not invent one; but the method for
deriving a local threshold is available and it does not require new data.

The threshold is set by two things: the consequence tier of
Section~\ref{sec:tiers}, and what control sits downstream of the assertion. Those two
bound the answer from opposite directions.

Where a competent reviewer examines every assertion before it is relied upon, a high
ungrounded rate is tolerable and the metric is not a safety measure at all---it is a
\emph{workload} measure, predicting how much reviewer time the system will consume and
whether the review is a rubber stamp. A system whose ungrounded rate exceeds what
reviewers can actually check has silently converted a control into a formality, and the
grounding rate is how that is detected.

Where the assertion is relied upon directly, the tolerable rate is bounded by something
the institution already knows. Every regulated firm has an accepted error rate for the
equivalent manual process: a sampling standard for file reviews, a tolerance for
unevidenced decisions found in quality assurance, an internal audit threshold that
triggers a finding. Those numbers exist, they were negotiated with a regulator or set by
policy, and they are the right anchor---because the question a supervisor will ask is not
whether the system is good in the abstract but whether it is at least as well evidenced
as the process it replaced. Anchoring to the incumbent standard also makes the comparison
honest in the direction that usually favors the system: manual processes frequently
cannot produce the evidence trail this book requires either.

Two cautions. Do not set a threshold on pooled accuracy and infer a grounding threshold
from it, which Corollary~\ref{cor:no-bound} forbids. And set the abstention and coverage
pair jointly, per Section~\ref{sec:abstention}, or the threshold will be met by a system
that declines to answer.

Collecting the four measurements with the two framing questions, a validation report
that supports sign-off contains, at minimum:

\begin{table}[ht]
  \centering
  \small
  \begin{tabular}{@{}>{\raggedright\arraybackslash}p{4.2cm}>{\raggedright\arraybackslash}p{4.0cm}>{\raggedright\arraybackslash}p{5.4cm}@{}}
    \toprule
    Statistic & Answers & Caveat to state \\
    \midrule
    Ungrounded assertion rate
      & how often it asserts without support
      & retrieval-grounded or parametric; rater agreement \\
    \addlinespace
    Misapplied rate
      & how often support is inapplicable at $c$
      & which components of $c$ were checked \\
    \addlinespace
    Error rate (grounded, false)
      & corpus quality
      & attributed to sources, not the system \\
    \addlinespace
    Abstention rate on held-out items
      & whether silence is available
      & corpus-controlled configuration only \\
    \addlinespace
    Form divergence
      & what a passing test is worth
      & probability-matching method \\
    \addlinespace
    Repeatability
      & whether verification persists
      & within- and cross-session \\
    \addlinespace
    Accuracy
      & performance on this sample
      & does not bound any of the above (Cor.~\ref{cor:no-bound}) \\
    \bottomrule
  \end{tabular}
  \caption{Accuracy stays on the report. It moves from the headline to a line item
    with a stated scope, which is what Corollary~\ref{cor:no-bound} licenses.}
  \label{tab:report}
\end{table}
